\documentclass[12pt]{article}
\usepackage[utf8]{inputenc}
\usepackage[spanish]{babel}
\usepackage[margin=2.5cm]{geometry}
\usepackage{listings}
\usepackage{xcolor}
\usepackage{parskip}
\usepackage{titlesec}
\usepackage{hyperref}

\definecolor{codebg}{RGB}{245,245,247}
\definecolor{codetext}{RGB}{30,30,30}

\lstset{
  backgroundcolor=\color{codebg},
  basicstyle=\ttfamily\small\color{codetext},
  breaklines=true,
  frame=single,
  rulecolor=\color{gray!40},
  framesep=6pt,
  xleftmargin=8pt,
  xrightmargin=8pt,
}

\titleformat{\section}{\large\bfseries}{}{0em}{}[\titlerule]
\titleformat{\subsection}{\normalsize\bfseries}{}{0em}{}

\hypersetup{colorlinks=true, linkcolor=black, urlcolor=blue}

\begin{document}

% ── Encabezado ─────────────────────────────────────────────────────────────────
\begin{center}
  {\LARGE \textbf{Linker}} \\[6pt]
  {\large Guía de Instalación y Despliegue} \\[4pt]
  \small Diego Cardenas \quad Samuel Albarracin \quad Carlos Barrero \quad Isaac Palomo
\end{center}

\vspace{6pt}
\hrule
\vspace{14pt}

% ── 1. Código fuente ───────────────────────────────────────────────────────────
\section{1. Acceder al código fuente}

Clonar el repositorio desde GitHub:

\begin{lstlisting}
git clone https://github.com/co-eiv-devsecops/linker3.git
\end{lstlisting}

% ── 2. Ejecutar localmente ────────────────────────────────────────────────────
\section{2. Ejecutar localmente}

Requiere \textbf{Node.js 22 o superior}. No hay dependencias externas que instalar.

\begin{lstlisting}
node server.js
\end{lstlisting}

Abrir en el navegador: \texttt{http://localhost:3000}

% ── 3. Despliegue en Oracle Cloud ────────────────────────────────────────────
\section{3. Despliegue en Oracle Cloud (Ubuntu)}

\textbf{URL de produccion:} \url{https://3.n-la-c.app} \\
\textbf{IP del servidor:} \texttt{10.0.69.233} \\
\textbf{Llave SSH:} \texttt{.ssh/linkervm-3.key}

\subsection{3.1 Conectarse a la máquina virtual}

\begin{lstlisting}
ssh -i .ssh/linkervm-3.key ubuntu@10.0.69.233
\end{lstlisting}

\subsection{3.2 Instalar Node.js (solo la primera vez)}

\begin{lstlisting}
sudo apt-get update
sudo apt-get install -y nodejs npm
\end{lstlisting}

Verificar versión:

\begin{lstlisting}
node --version
\end{lstlisting}

\subsection{3.3 Copiar archivos desde la máquina local}

Ejecutar desde tu computador (no desde la VM):

\begin{lstlisting}
scp -i .ssh/linkervm-3.key -r ~/ruta/Tarea1/ ubuntu@10.0.69.233:~/dummy/
\end{lstlisting}

\subsection{3.4 Configurar el servicio systemd}

Dentro de la VM, crear el archivo de servicio:

\begin{lstlisting}
sudo nano /etc/systemd/system/linker.service
\end{lstlisting}

Pegar el siguiente contenido:

\begin{lstlisting}
[Unit]
Description=Linker Node.js
After=network.target

[Service]
Type=simple
User=ubuntu
WorkingDirectory=/home/ubuntu/dummy
Environment=PORT=8080
ExecStart=/usr/bin/node /home/ubuntu/dummy/server.js
Restart=always
RestartSec=10
KillSignal=SIGINT
SyslogIdentifier=linker

[Install]
WantedBy=multi-user.target
\end{lstlisting}

Activar e iniciar:

\begin{lstlisting}
sudo systemctl daemon-reload
sudo systemctl enable linker
sudo systemctl start linker
\end{lstlisting}

\subsection{3.5 Verificar que está corriendo}

\begin{lstlisting}
sudo systemctl status linker
\end{lstlisting}

% ── 4. Actualizar la aplicación ───────────────────────────────────────────────
\section{4. Actualizar la aplicación}

\textbf{Paso 1.} Copiar los archivos modificados desde tu computador:

\begin{lstlisting}
scp -i .ssh/linkervm-3.key server.js ubuntu@10.0.69.233:~/dummy/server.js
scp -i .ssh/linkervm-3.key public/index.html ubuntu@10.0.69.233:~/dummy/public/index.html
\end{lstlisting}

\textbf{Paso 2.} Reiniciar el servicio en la VM:

\begin{lstlisting}
sudo systemctl restart linker
\end{lstlisting}

\textbf{Paso 3.} Verificar logs si algo falla:

\begin{lstlisting}
journalctl -u linker -f
\end{lstlisting}

% ── 5. Estructura de archivos en la VM ───────────────────────────────────────
\section{5. Estructura de archivos en la VM}

\begin{lstlisting}
ubuntu@linker3:~/dummy/
|-- server.js
|-- package.json
|-- linker.db
+-- public/
    +-- index.html
\end{lstlisting}

% ── 6. Script de despliegue (Bono) ───────────────────────────────────────────
\section{6. Script de despliegue automatico}

El archivo \texttt{deploy.sh} copia los archivos al servidor y reinicia el servicio en un solo comando.

\subsection{Requisitos previos}

Tener la llave SSH en \texttt{.ssh/linkervm-3.key} en tu computador local.

\subsection{Invocar el script}

Desde la raiz del proyecto en tu computador local:

\begin{lstlisting}
bash deploy.sh
\end{lstlisting}

El script realiza automaticamente:

\begin{enumerate}
  \item Copia \texttt{server.js}, \texttt{public/index.html} y \texttt{package.json} a \texttt{/home/ubuntu/dummy/} via SCP.
  \item Reinicia el servicio \texttt{linker} con \texttt{systemctl restart}.
\end{enumerate}

No es necesario conectarse manualmente a la VM para una actualizacion normal.

\end{document}
